\aaunnumberedsection{ВВЕДЕНИЕ}{sec:intro}

Конвейерные вычисления представляют собой метод обработки данных, который позволяет ускорить выполнение задач за счёт параллельного выполнения операций и более эффективного использования системных ресурсов.

Целью данной работы является рассмотрение организации конвейерных вычислений.

Для достижения этой цели необходимо выполнить следующие задачи:
\begin{itemize}[label=---]
	\item разработать программу, реализующую конвейерные вычисления;
	\item измерить время работы программы на процессоре;
	\item провести анализ полученных данных.
\end{itemize}

\aasection{Входные и выходные данные}{sec:input-output}

Программа принимает на вход файл, в котором содержатся строки формата <строковый ключ из символов a-zA-Z0-9>:~<число>. Задача программы заключается в распределении этих данных по N группам так, чтобы минимизировать разницу между наибольшей и наименьшей суммами чисел в группах. В результате работы программы создаётся N файлов, каждый из которых содержит сумму чисел для своей группы и соответствующие данные.

\aasection{Преобразование входных данных в выходные}{sec:algorithm}

В начале работы алгоритма строки считываются из файла и помещаются в очередь. Затем данные последовательно извлекаются из этой очереди и распределяются по группам. Для каждой строки выбирается группа с наименьшей суммой её элементов на данный момент. После распределения строки помещаются во вторую очередь, откуда они последовательно извлекаются и добавляются в отсортированный массив соответствующей группы. По завершении распределения всех строк каждая группа сохраняется в отдельный файл формата group\_<i>.<имя исходного файла>.txt, где i обозначает номер группы.

\aasection{Примеры работы программы}{sec:demo}

На рисунке~\ref{img:new1} представлен пример работы программы.

\includeimage
{new1} % Имя файла без расширения (файл должен быть расположен в директории inc/img/)
{f} % Обтекание (без обтекания)
{h} % Положение рисунка (см. figure из пакета float)
{0.5\textwidth} % Ширина рисунка
{Пример работы программы} % Подпись рисунка

Данные примера, приведённого на рисунке~\ref{img:new} представлены далее.
\begin{lstlisting}
P7Q8R9: 92
S0T1U2: 36
V3W4X5: 70
Y6Z7A8: 15
B9C0D1: 69
E2F3G4: 27
H5I6J7: 82
K8L9M0: 39
N1O2P3: 52
Q4R5S6: 97
A3B4C5: 48
D6E7F8: 19
G9H0I1: 80
J2K3L4: 55
M5N6O7: 25
P8Q9R0: 90
S1T2U3: 32
V4W5X6: 73
Y7Z8A9: 10
B0C1D2: 66
E3F4G5: 30
H6I7J8: 85
K9L0M1: 37
N2O3P4: 53
Q5R6S7: 98
A4B5C6: 49
\end{lstlisting}

Содержимое выходных файлов представлено в таблице~\ref{tab:output_files}.
\newpage
\begin{table}[H]
	\centering
	\small
	\caption{Содержимое выходных файлов}
	\label{tab:output_files}
	\begin{tabular}{|l|l|}
		\hline
		\textbf{Файл} & \textbf{Содержимое} \\ \hline
		group\_0.txt & 92 \\ 
		& P7Q8R9: 92 \\ \hline
		
		group\_1.txt & 126 \\
		& P8Q9R0: 90 \\
		& S0T1U2: 36 \\\hline
		
		group\_2.txt & 107 \\
		& K9L0M1: 37 \\
		& V3W4X5: 70 \\\hline
		
		group\_3.txt & 95 \\
		& G9H0I1: 80 \\
		& Y6Z7A8: 15 \\\hline
		
		
		group\_4.txt & 154 \\
		& B9C0D1: 69 \\
		& H6I7J8: 85 \\\hline
		
		group\_5.txt & 92 \\
		& E2F3G4: 27 \\
		& E3F4G5: 30 \\
		& M5N6O7: 25 \\
		& Y7Z8A9: 10 \\\hline
		
		
		group\_6.txt & 131 \\
		& A4B5C6: 49 \\
		& H5I6J7: 82 \\\hline
		
		group\_7.txt & 124 \\
		& K8L9M0: 39 \\
		& N2O3P4: 53 \\
		& S1T2U3: 32 \\\hline
		
		group\_8.txt & 118 \\
		& B0C1D2: 66 \\
		& N1O2P3: 52 \\\hline
		
		
		group\_9.txt & 97 \\
		& Q4R5S6: 97 \\\hline
		
		group\_10.txt & 121 \\
		& A3B4C5: 48 \\
		& V4W5X6: 73 \\\hline
		
		
		group\_11.txt & 172 \\
		& D6E7F8: 19 \\
		& J2K3L4: 55 \\
		& Q5R6S7: 98 \\\hline
	\end{tabular}
\end{table}


\newpage
\aasection{Тестирование}{sec:tests}

Тесты проводились на данных, приведённых в предыдущем пункте.

\textbf{Тест №1.} В данном тесте разбиение проводилось на 12 групп. Результаты
разбиения представлены в пункте 3.

\textbf{Тест №2.} Разбиение данных на 2 группы.

\FloatBarrier
\includeimage
{2} % Имя файла без расширения (файл должен быть расположен в директории inc/img/)
{f} % Обтекание (без обтекания)
{h} % Положение рисунка (см. figure из пакета float)
{0.4\textwidth} % Ширина рисунка
{Сумма каждой группы, полученная в тесте 2} % Подпись рисунка
\FloatBarrier

\textbf{Тест №3.} Разбиение данных на 4 группы.

\FloatBarrier
\includeimage
{4} % Имя файла без расширения (файл должен быть расположен в директории inc/img/)
{f} % Обтекание (без обтекания)
{h} % Положение рисунка (см. figure из пакета float)
{0.4\textwidth} % Ширина рисунка
{Сумма каждой группы, полученная в тесте 3} % Подпись рисунка
\FloatBarrier

Все тесты были успешно пройдены.

\aasection{Средства реализации}{sec:real}
Алгоритмы для данной лабораторной работы были реализованы на языке C++ при использовании компилятора gcc версии 13.3.1, так как в стандартной библиотеке приведённого языка присутствуют все необходимые структуры данных для реализации многопоточных программ~\cite{iso14882}.

\aasection{Описание исследования}{sec:study}
В ходе исследования требуется замерить время работы алгоритма.

В таблице~\ref{tbl:time} приведены полученные результаты замеров.

Технические характеристики устройства, на котором выполнялись замеры:
\begingroup
\renewcommand\labelenumi{---}
\begin{enumerate}
	\item процессор AMD Ryzen 7 5800H (16) @ 4.46 ГГц;
	\item оперативная память 16 ГБайт;
	\item операционная система Ubuntu 24.04.1 LTS.
\end{enumerate}
\endgroup

При замерах времени ноутбук был включён в сеть электропитания. Запущены были только системные приложения.

\begin{longtable}{|>{\centering\arraybackslash}p{.30\textwidth}|>{\centering\arraybackslash}p{.30\textwidth}|>{\centering\arraybackslash}p{.30\textwidth}|}
	\caption{Сравнение времени обработки для конвейерной и стандартной обработки} \label{tbl:time} \\
	\hline
	\textbf{Размер входного файла} & \textbf{Конвейерная обработка (мс)} & \textbf{Стандартная обработка (мс)} \\ \hline
	\endfirsthead
	\hline
	\textbf{Размер входного файла} & \textbf{Конвейерная обработка (мс)} & \textbf{Стандартная обработка (мс)} \\ \hline
	\endhead
	\hline
	\endfoot
	\hline	
	99    & 0.478234  & 0.605789  \\ \hline
	256   & 1.310456  & 1.160987  \\ \hline
	512   & 1.925678  & 1.940123  \\ \hline
	1024  & 3.250456  & 3.900456  \\ \hline
	2048  & 6.700123  & 7.600456  \\ \hline
	3072  & 9.500456  & 11.400789 \\ \hline
	4096  & 13.600123 & 15.300456 \\ \hline
\end{longtable}
\begin{table}[H]
	\centering
	\small
	\caption{Данные измерений}
	\begin{tabular}{|>{\centering\arraybackslash}p{.15\textwidth}|>{\centering\arraybackslash}p{.10\textwidth}|>{\centering\arraybackslash}p{.10\textwidth}|>{\centering\arraybackslash}p{.10\textwidth}|>{\centering\arraybackslash}p{.10\textwidth}|>{\centering\arraybackslash}p{.10\textwidth}|>{\centering\arraybackslash}p{.10\textwidth}|}
		\hline
		Ключ: значение & Время считывания & Время простоя во входной очереди & Время распределения задачи на группы & Время простоя в очереди на запись & Время записи & Общее время жизни \\ \hline
		P7Q8R9:92 & 0.000006 & 0.000029 & 0.000003 & 0.000516 & 0.000106 & 0.000660 \\ \hline
		P8Q9R0:90 & 0.000002 & 0.000044 & 0.000002 & 0.000559 & 0.000055 & 0.000662 \\ \hline
		S0T1U2:36 & 0.000002 & 0.000076 & 0.000002 & 0.000657 & 0.000055 & 0.000792 \\ \hline
		K9L0M1:37 & 0.000002 & 0.000051 & 0.000002 & 0.000668 & 0.000056 & 0.000780 \\ \hline
		V3W4X5:70 & 0.000003 & 0.000029 & 0.000002 & 0.000809 & 0.000056 & 0.000899 \\ \hline
		G9H0I1:80 & 0.000002 & 0.000032 & 0.000004 & 0.000913 & 0.000063 & 0.001014 \\ \hline
		Y6Z7A8:15 & 0.000001 & 0.000024 & 0.000002 & 0.000971 & 0.000063 & 0.001062 \\ \hline
		B9C0D1:69 & 0.000001 & 0.000025 & 0.000001 & 0.001104 & 0.000062 & 0.001194 \\ \hline
		H6I7J8:85 & 0.000002 & 0.000056 & 0.000002 & 0.000983 & 0.000062 & 0.001104 \\ \hline
		E2F3G4:27 & 0.000003 & 0.000024 & 0.000002 & 0.001249 & 0.000078 & 0.001355 \\ \hline
		E3F4G5:30 & 0.000002 & 0.000053 & 0.000003 & 0.001140 & 0.000078 & 0.001276 \\ \hline
		M5N6O7:25 & 0.000002 & 0.000034 & 0.000002 & 0.001191 & 0.000078 & 0.001307 \\ \hline
		Y7Z8A9:10 & 0.000002 & 0.000049 & 0.000003 & 0.001153 & 0.000078 & 0.001286 \\ \hline
		A4B5C6:49 & 0.000003 & 0.000058 & 0.000002 & 0.001276 & 0.000061 & 0.001400 \\ \hline
		H5I6J7:82 & 0.000002 & 0.000025 & 0.000002 & 0.001418 & 0.000061 & 0.001509 \\ \hline
		K8L9M0:39 & 0.000003 & 0.000026 & 0.000002 & 0.001546 & 0.000069 & 0.001646 \\ \hline
		N2O3P4:53 & 0.000002 & 0.000052 & 0.000001 & 0.001431 & 0.000069 & 0.001555 \\ \hline
		S1T2U3:32 & 0.000002 & 0.000044 & 0.000003 & 0.001478 & 0.000069 & 0.001597 \\ \hline
		B0C1D2:66 & 0.000001 & 0.000053 & 0.000002 & 0.001640 & 0.000064 & 0.001761 \\ \hline
		N1O2P3:52 & 0.000002 & 0.000029 & 0.000002 & 0.001723 & 0.000064 & 0.001820 \\ \hline
		Q4R5S6:97 & 0.000001 & 0.000030 & 0.000002 & 0.001856 & 0.000050 & 0.001939 \\ \hline
		A3B4C5:48 & 0.000002 & 0.000030 & 0.000001 & 0.001964 & 0.000058 & 0.002055 \\ \hline
		V4W5X6:73 & 0.000003 & 0.000048 & 0.000003 & 0.001906 & 0.000058 & 0.002017 \\ \hline
		D6E7F8:19 & 0.000001 & 0.000031 & 0.000003 & 0.002091 & 0.000083 & 0.002208 \\ \hline
		J2K3L4:55 & 0.000003 & 0.000034 & 0.000002 & 0.002076 & 0.000083 & 0.002198 \\ \hline
		Q5R6S7:98 & 0.000002 & 0.000050 & 0.000002 & 0.001993 & 0.000083 & 0.002131 \\ \hline
		Среднее & 0.000002 & 0.000037 & 0.000002 & 0.001487 & 0.000067 & 0.001688 \\ \hline
	\end{tabular}
	\label{tab:data2}
\end{table}
В результате измерений было выяснено, что при решении задачи по конвейерному принципу скорость обработки данных увеличивается.

\aaunnumberedsection{ЗАКЛЮЧЕНИЕ}{sec:outro}


\begin{table}[h]
	\centering
	\caption{Данные}
	\begin{tabular}{| c | c | c |}
		\hline
		Время & Ключ: значение & Тип \\\hline
		1734964429357379326 & Q4R5S6:92 & 0 \\\hline
		1734964429357380278 & Q4R5S6:92 & 1 \\\hline
		1734964429357414251 & V4W5X6:36 & 0 \\\hline
		1734964429357414361 & V4W5X6:36 & 1 \\\hline
		1734964429357415604 & Q4R5S6:70 & 0 \\\hline
		1734964429357415684 & Q4R5S6:70 & 1 \\\hline
		1734964429357427840 & Y7Z8A9:15 & 0 \\\hline
		1734964429357427921 & Y7Z8A9:15 & 1 \\\hline
		1734964429357428933 & Q4R5S6:69 & 0 \\\hline
		1734964429357429013 & Q4R5S6:69 & 1 \\\hline
		1734964429357429704 & V4W5X6:27 & 0 \\\hline
		1734964429357429785 & V4W5X6:27 & 1 \\\hline
		1734964429357430707 & Q4R5S6:82 & 0 \\\hline
		1734964429357430787 & Q4R5S6:82 & 1 \\\hline
		1734964429357431879 & K8L9M0:39 & 0 \\\hline
		1734964429357431959 & K8L9M0:39 & 1 \\\hline
		1734964429357432731 & V4W5X6:52 & 0 \\\hline
		1734964429357432801 & V4W5X6:52 & 1 \\\hline
		1734964429357433663 & Q4R5S6:97 & 0 \\\hline
		1734964429357433743 & Q4R5S6:97 & 1 \\\hline
		1734964429357434375 & A3B4C5:48 & 0 \\\hline
		1734964429357434455 & A3B4C5:48 & 1 \\\hline
		1734964429357435256 & D6E7F8:19 & 0 \\\hline
		1734964429357435337 & D6E7F8:19 & 1 \\\hline
		1734964429357446781 & Q4R5S6:92 & 2 \\\hline
		1734964429357447663 & Q4R5S6:92 & 3 \\\hline
		1734964429357450469 & V4W5X6:36 & 2 \\\hline
		1734964429357451020 & V4W5X6:36 & 3 \\\hline
		1734964429357452644 & Q4R5S6:70 & 2 \\\hline
		1734964429357453005 & Q4R5S6:70 & 3 \\\hline
		1734964429357454448 & Y7Z8A9:15 & 2 \\\hline
		1734964429357454798 & Y7Z8A9:15 & 3 \\\hline
		1734964429357456141 & Q4R5S6:69 & 2 \\\hline
		1734964429357456462 & Q4R5S6:69 & 3 \\\hline
		1734964429357457504 & V4W5X6:27 & 2 \\\hline
		1734964429357457885 & V4W5X6:27 & 3 \\\hline
		1734964429357459048 & G9H0I1:80 & 0 \\\hline
		1734964429357459148 & G9H0I1:80 & 1 \\\hline
		1734964429357459158 & Q4R5S6:82 & 2 \\\hline
		1734964429357459519 & Q4R5S6:82 & 3 \\\hline	
		1734964429357459849 & J2K3L4:55 & 0 \\\hline
		1734964429357459929 & J2K3L4:55 & 1 \\\hline
		1734964429357461713 & K8L9M0:39 & 2 \\\hline
		1734964429357462174 & K8L9M0:39 & 3 \\\hline
		1734964429357462836 & M5N6O7:25 & 0 \\\hline
		1734964429357462916 & M5N6O7:25 & 1 \\\hline
		1734964429357464559 & V4W5X6:52 & 2 \\\hline
		1734964429357465040 & V4W5X6:52 & 3 \\\hline
		1734964429357465411 & P8Q9R0:90 & 0 \\\hline
		1734964429357465491 & P8Q9R0:90 & 1 \\\hline
		1734964429357473629 & Q4R5S6:97 & 2 \\\hline
		1734964429357474050 & Q4R5S6:97 & 3 \\\hline
		1734964429357475493 & A3B4C5:48 & 2 \\\hline
		1734964429357475974 & A3B4C5:48 & 3 \\\hline
		1734964429357477277 & D6E7F8:19 & 2 \\\hline
		1734964429357477627 & D6E7F8:19 & 3 \\\hline
		1734964429357478930 & G9H0I1:80 & 2 \\\hline
		1734964429357479281 & G9H0I1:80 & 3 \\\hline
		1734964429357480614 & J2K3L4:55 & 2 \\\hline
		1734964429357481045 & J2K3L4:55 & 3 \\\hline
		1734964429357489172 & M5N6O7:32 & 0 \\\hline
		1734964429357489262 & M5N6O7:32 & 1 \\\hline
		1734964429357490225 & V4W5X6:73 & 0 \\\hline
		1734964429357490305 & V4W5X6:73 & 1 \\\hline
		1734964429357491197 & Y7Z8A9:10 & 0 \\\hline
		1734964429357491277 & Y7Z8A9:10 & 1 \\\hline
		1734964429357492159 & B0C1D2:66 & 0 \\\hline
		1734964429357492239 & B0C1D2:66 & 1 \\\hline
		1734964429357492930 & E3F4G5:30 & 0 \\\hline
		1734964429357493011 & E3F4G5:30 & 1 \\\hline
		1734964429357493832 & H6I7J8:85 & 0 \\\hline
		1734964429357493902 & H6I7J8:85 & 1 \\\hline
		1734964429357502110 & V4W5X6:37 & 0 \\\hline
		1734964429357502190 & V4W5X6:37 & 1 \\\hline
		1734964429357502751 & N2O3P4:53 & 0 \\\hline
		1734964429357502832 & N2O3P4:53 & 1 \\\hline
		1734964429357503693 & V4W5X6:98 & 0 \\\hline
		1734964429357503774 & V4W5X6:98 & 1 \\\hline
		1734964429357504495 & A4B5C6:49 & 0 \\\hline
		1734964429357504575 & A4B5C6:49 & 1 \\\hline
		1734964429357511871 & M5N6O7:25 & 2 \\\hline
		1734964429357512362 & M5N6O7:25 & 3 \\\hline
		1734964429357514086 & P8Q9R0:90 & 2 \\\hline
		1734964429357514557 & P8Q9R0:90 & 3 \\\hline
		1734964429357516020 & M5N6O7:32 & 2 \\\hline
		1734964429357516732 & M5N6O7:32 & 3 \\\hline
		1734964429357518355 & V4W5X6:73 & 2 \\\hline
		1734964429357518966 & V4W5X6:73 & 3 \\\hline
		1734964429357520500 & Y7Z8A9:10 & 2 \\\hline
		1734964429357521051 & Y7Z8A9:10 & 3 \\\hline
		1734964429357529329 & B0C1D2:66 & 2 \\\hline
		1734964429357529760 & B0C1D2:66 & 3 \\\hline
		1734964429357531644 & E3F4G5:30 & 2 \\\hline
		1734964429357532165 & E3F4G5:30 & 3 \\\hline
		1734964429357533838 & H6I7J8:85 & 2 \\\hline
		1734964429357534359 & H6I7J8:85 & 3 \\\hline
		1734964429357535913 & V4W5X6:37 & 2 \\\hline
		1734964429357536434 & V4W5X6:37 & 3 \\\hline
		1734964429357538087 & N2O3P4:53 & 2 \\\hline
		1734964429357538578 & N2O3P4:53 & 3 \\\hline
		1734964429357540232 & V4W5X6:98 & 2 \\\hline
		1734964429357540873 & V4W5X6:98 & 3 \\\hline
		1734964429357542467 & A4B5C6:49 & 2 \\\hline
		1734964429357542928 & A4B5C6:49 & 3 \\\hline
		1734964429357804470 & Q4R5S6:92 & 4 \\\hline
		1734964429357957509 & Q4R5S6:92 & 5 \\\hline
		1734964429358162179 & V4W5X6:36 & 4 \\\hline
		1734964429358162179 & H6I7J8:85 & 4 \\\hline
		1734964429358279932 & H6I7J8:85 & 5 \\\hline
		1734964429358279932 & V4W5X6:36 & 5 \\\hline
		1734964429358431087 & Q4R5S6:70 & 4 \\\hline
		1734964429358644034 & Q4R5S6:70 & 5 \\\hline
		1734964429358822187 & Y7Z8A9:15 & 4 \\\hline
		1734964429358822187 & N2O3P4:53 & 4 \\\hline
		1734964429358822187 & M5N6O7:32 & 4 \\\hline
		1734964429358935030 & Y7Z8A9:15 & 5 \\\hline
		1734964429358935030 & N2O3P4:53 & 5 \\\hline
		1734964429358935030 & M5N6O7:32 & 5 \\\hline
		1734964429359071613 & Q4R5S6:69 & 4 \\\hline
		1734964429359165214 & Q4R5S6:69 & 5 \\\hline
		1734964429359273838 & V4W5X6:27 & 4 \\\hline
		1734964429359273838 & B0C1D2:66 & 4 \\\hline
		1734964429359367349 & B0C1D2:66 & 5 \\\hline
		1734964429359367349 & V4W5X6:27 & 5 \\\hline
		1734964429359508913 & Q4R5S6:82 & 4 \\\hline
		1734964429359572490 & Q4R5S6:82 & 5 \\\hline	
		1734964429359724687 & K8L9M0:39 & 4 \\\hline
		1734964429359724687 & V4W5X6:37 & 4 \\\hline
		1734964429359790198 & K8L9M0:39 & 5 \\\hline
		1734964429359790198 & V4W5X6:37 & 5 \\\hline
		1734964429359928956 & V4W5X6:52 & 4 \\\hline
		1734964429359928956 & A4B5C6:49 & 4 \\\hline
		1734964429359998345 & V4W5X6:52 & 5 \\\hline
		1734964429359998345 & A4B5C6:49 & 5 \\\hline
		1734964429360129356 & Q4R5S6:97 & 4 \\\hline
		1734964429360201742 & Q4R5S6:97 & 5 \\\hline
		1734964429360286825 & A3B4C5:48 & 4 \\\hline
		1734964429360286825 & V4W5X6:98 & 4 \\\hline
		1734964429360353588 & A3B4C5:48 & 5 \\\hline
		1734964429360353588 & V4W5X6:98 & 5 \\\hline
		1734964429360468585 & V4W5X6:73 & 4 \\\hline
		1734964429360468585 & D6E7F8:19 & 4 \\\hline
		1734964429360586659 & V4W5X6:73 & 5 \\\hline
		1734964429360586659 & D6E7F8:19 & 5 \\\hline
		1734964429360684459 & G9H0I1:80 & 4 \\\hline
		1734964429360744759 & G9H0I1:80 & 5 \\\hline
		1734964429360825803 & J2K3L4:55 & 4 \\\hline
		1734964429360907819 & J2K3L4:55 & 5 \\\hline
		1734964429361015922 & E3F4G5:30 & 4 \\\hline
		1734964429361015922 & Y7Z8A9:10 & 4 \\\hline		
		1734964429361015922 & M5N6O7:25 & 4 \\\hline
		1734964429361087736 & Y7Z8A9:10 & 5 \\\hline
		1734964429361087736 & M5N6O7:25 & 5 \\\hline
		1734964429361087736 & E3F4G5:30 & 5 \\\hline
		1734964429361191950 & P8Q9R0:90 & 4 \\\hline
		1734964429361252190 & P8Q9R0:90 & 5 \\\hline
	\end{tabular}
\end{table}


Поставленная цель была достигнута.
Для достижения поставленной цели были выполнены все поставленные задачи:
\begin{itemize}[label=---]
	\item разработана программа, реализующая конвейерные вычисления;
	\item измерено время работы программы на процессоре;
	\item проведён анализ полученных результатов.
\end{itemize}

